% ============================================================================
%  Chapter 3 -- History of Computing                             (Lecture 03)
% ============================================================================
\chapter{History of Computing}
\label{ch:history}

\epigraphline{\texttt{printf("Hello, world!");}}
  {Brian Kernighan and Dennis Ritchie, \emph{The C Programming Language}, 1978}

\begin{chapgoals}
By the end of this chapter you will be able to
\begin{itemize}
\item define informatics and say why it is not the science of computers;
\item place the key machines and people on a timeline, from the abacus to the
  stored-program computer;
\item name the four units of the von Neumann architecture and find them in the
  machine programmed in Chapter~\ref{ch:firstcontact};
\item explain what RISC means and why \riscv is an open instruction set;
\item describe where informatics sits between science, engineering and the economy.
\end{itemize}
\end{chapgoals}
% Generated by tools/textbook/export_listings.js from tools/decks/mastery.js and tools/concepts/registry.json: do not edit.
\begin{chapstatus}{Learning status of this chapter}
\masterychip{now}{UNDERSTAND NOW}\enspace what informatics is; the von Neumann architecture; RISC and an open ISA.\par\smallskip
\masterychip{use}{USE NOW}\enspace the complete program, the PRG32 tools (used, not studied, in this course).\par\smallskip
\masterychip{later}{UNDERSTAND LATER}\enspace Boolean algebra (Chapter~\ref{ch:boole1}).
\end{chapstatus}
\begin{chapstatus}{Roadmap of this chapter}
By the end of the chapter you can:
\begin{enumerate}[nosep,leftmargin=2em]
\item define informatics and explain why it is not only about computers
\item name the units of the von Neumann architecture and its stored-program idea
\item say what distinguishes RISC from CISC and why RISC-V is open
\end{enumerate}
It builds on: only the setup of Chapter~\ref{ch:setup}. Next: Chapter~\ref{ch:information}.
\end{chapstatus}


The reader has already run a \riscv program. This chapter tells where that
machine comes from. It is not a chronicle: of the thousands of machines and
people that the history of computing could list, it keeps those that added an
idea which is still inside the ESP32-C6 on the breadboard. The thread to
follow is that, at every stage, somebody found a new way of doing one of two
things: \emph{representing} information in a physical object, or
\emph{prescribing} a sequence of operations to be carried out without further
thought.

% ----------------------------------------------------------------------------
\section{What is informatics?}

\begin{definition}{Informatics}
Informatics is the applied science that deals with the treatment of
information by means of automated procedures.
\end{definition}

The Italian word \emph{informatica} comes from the French \emph{informatique},
a contraction of \emph{information automatique}. English says \emph{computer
science} and, increasingly, \emph{informatics}. The tool of the discipline
is called \emph{elaboratore} or \emph{calcolatore} in Italian,
\emph{ordinateur} in French, \emph{computer} in English. The English name of
the discipline has the defect of putting the tool in the place of the subject,
a defect that a sentence attributed to Edsger W.\ Dijkstra corrects:

\begin{quote}\itshape
Computer science is no more about computers than astronomy is about telescopes.
\end{quote}

The machine is the instrument, not the subject. In this book the instrument
is studied closely precisely so that it stops being a mystery.

The discipline has two souls. \emph{Theoretical informatics} studies the
foundations of information and computation at the logical level: what can be
computed, and at what cost. Its central object is the \term{algorithm}, a
finite set of operations that describes the transformation of a datum $A$ into
a result $B$. \emph{Applied informatics} studies the enabling technologies
and the information systems built with them. It is pervasive, a driver of
economic development and of the whole sector of information and communication
technologies, a STEM discipline at the heart of an industrial revolution; and
it has a social cost that must be managed, the \term{digital divide} between
those who have access to these technologies and those who do not.

% ----------------------------------------------------------------------------
\section{Before electricity}

For three thousand years people have built machines to help them calculate.

\subsection{Origins: the abacus and the Antikythera mechanism}

Babylonians, Chinese, Greeks and Romans all used counting frames
(Figure~\ref{fig:hist:origins}). In an abacus a number \emph{is} a position
of beads, and an addition \emph{is} a movement of the fingers: the device is a
physical representation of data and of operations. The Antikythera mechanism,
recovered from a shipwreck and built between the second and the first century
before Christ (most estimates fall between about 205 and 60~BC), is the oldest
known geared calculator: its dials computed astronomical positions and
eclipses, and in it a number is the angle of a gear.

\begin{figure}[htb]
\centering
\begin{subfigure}[b]{0.24\textwidth}\includegraphics[width=\textwidth,height=32mm,keepaspectratio]{abacus_roman.jpg}\caption{Roman hand abacus}\end{subfigure}\hfill
\begin{subfigure}[b]{0.24\textwidth}\includegraphics[width=\textwidth,height=32mm,keepaspectratio]{abacus_chinese.jpg}\caption{Chinese \emph{suanpan}}\end{subfigure}\hfill
\begin{subfigure}[b]{0.24\textwidth}\includegraphics[width=\textwidth,height=32mm,keepaspectratio]{antikythera.jpg}\caption{Antikythera mechanism}\end{subfigure}\hfill
\begin{subfigure}[b]{0.24\textwidth}\includegraphics[width=\textwidth,height=32mm,keepaspectratio]{antikythera_plate.jpg}\caption{A modern model of its dials}\end{subfigure}
\caption{Origins: a number becomes a position of beads or the angle of a gear.}
\label{fig:hist:origins}
\end{figure}

\begin{prgconcept}
Representation comes before calculation. Before anything can be computed,
something physical must stand for a number.
\end{prgconcept}

\subsection{The medieval Islamic world}

Between the ninth and the thirteenth century the scholars of the Islamic world
preserved the science of antiquity and added to it three things that concern
this story. The first is a family of astronomical instruments, the mechanical
astrolabe (Figure~\ref{fig:hist:astrolabe}) and the torquetum, which are
analogue computers for the positions of the stars. The second is
cryptanalysis: in the ninth century al-Kindi described frequency analysis, the
method of breaking a cipher by counting how often each symbol occurs. The
third is the programmable machine: the automatic flute player of the Ban\=u
M\=us\=a brothers, and the weight-driven clocks and automata of al-Jazari
(1206), whose behaviour could be changed by rearranging pegs.

The very word \emph{algorithm} comes from the name of the mathematician
al-Khw\=arizm\=\i, as \emph{algebra} comes from the title of his book.

\begin{figure}[htb]
\centering
\begin{subfigure}[b]{0.3\textwidth}\centering\includegraphics[width=\textwidth,height=42mm,keepaspectratio]{astrolabe.jpg}\caption{A brass astrolabe}\label{fig:hist:astrolabe}\end{subfigure}\hfill
\begin{subfigure}[b]{0.3\textwidth}\centering\includegraphics[width=\textwidth,height=42mm,keepaspectratio]{pascaline.jpg}\caption{A Pascaline}\label{fig:hist:pascaline}\end{subfigure}\hfill
\begin{subfigure}[b]{0.3\textwidth}\centering\includegraphics[width=\textwidth,height=42mm,keepaspectratio]{leibniz_reckoner.jpg}\caption{Leibniz's reckoner}\end{subfigure}
\caption{Instruments and calculators before 1700.}
\end{figure}

\subsection{From 1600 to 1800: mechanical calculators and logic}

The seventeenth century built the first mechanical calculators. Napier's
tables of logarithms (1614) turned multiplications into additions. Blaise
Pascal's \emph{Pascaline} (1642) added and subtracted, and its real invention
was the automatic \term{carry}: each wheel is one decimal digit
(Figure~\ref{fig:hist:pascaline}), and a wheel that passes from 9 to 0 pushes
its neighbour forward by one step, exactly as the full adder of
Chapter~\ref{ch:combinational} will do with bits. Leibniz's stepped reckoner
(1673) performed the four operations.

The same Leibniz described binary arithmetic in 1703, and a century and a half
later George Boole, in \emph{The Laws of Thought} (1854), showed that logic
can be treated as an algebra \citep{boole1854}: the subject of Chapters
\ref{ch:boole1} and~\ref{ch:boole2}.

Meanwhile a machine that had nothing to do with numbers introduced the stored
program. In Jacquard's loom (1804) the pattern of the cloth is recorded on a
chain of punched cards: change the cards, and the same loom weaves a different
pattern.

\subsection{The 1800s: Charles Babbage and Ada Lovelace}

Charles Babbage put the two lines together
(Figure~\ref{fig:hist:babbage}). His \emph{Difference Engine} computed
mathematical tables automatically. His \emph{Analytical Engine}, designed but
never completed, was a general-purpose computer in everything but technology:
it had a \emph{mill}, the arithmetic unit; a control unit; a \emph{store}, the
memory; and an output device. Programs and data were to be supplied on
punched cards, borrowed from the Jacquard loom.

Ada Lovelace devised a method for programming the Analytical Engine. Her
notes of 1843 contain what is regarded as the first published program, for
computing the Bernoulli numbers, and she is regarded as the first programmer
in history.

\begin{figure}[htb]
\centering
\begin{subfigure}[b]{0.32\textwidth}\centering\includegraphics[width=\textwidth,height=40mm,keepaspectratio]{difference_engine.jpg}\caption{Difference Engine}\end{subfigure}\hfill
\begin{subfigure}[b]{0.32\textwidth}\centering\includegraphics[width=\textwidth,height=40mm,keepaspectratio]{analytical_engine.jpg}\caption{Analytical Engine}\end{subfigure}\hfill
\begin{subfigure}[b]{0.32\textwidth}\centering\includegraphics[width=\textwidth,height=40mm,keepaspectratio]{punched_cards.jpg}\caption{Punched cards}\end{subfigure}
\caption{The machines of Charles Babbage.}
\label{fig:hist:babbage}
\end{figure}

% ----------------------------------------------------------------------------
\section{The twentieth century}

The twentieth century produced the idea of a universal machine, and then the
machines that made it real. The theory came first.

\subsection{Turing, Shannon and Nakajima}

In 1936 Alan Turing defined an abstract universal machine, and with it the
notion of what can be computed \citep{turing1936}. During the war he worked
at Bletchley Park on breaking the German Enigma cipher
(Figure~\ref{fig:hist:enigma}).

In the same years Akira Nakajima in Japan and Claude Shannon in the United
States founded the theory of switching circuits: Boolean algebra describes
circuits made of switches \citep{shannon1938}. It is the principle on which
all digital design rests, and the bridge between the logic of Chapters
\ref{ch:boole1}--\ref{ch:boole2} and the circuits of Chapters
\ref{ch:electronics}--\ref{ch:combinational}.

In Germany Konrad Zuse built the Z3 (1941), the first working computer
controlled by a program, which was shown in 1998 to be Turing-complete; its
successor, the Z4, was one of the first computers to be sold commercially.

\begin{figure}[htb]
\centering
\begin{subfigure}[b]{0.32\textwidth}\centering\includegraphics[width=\textwidth,height=38mm,keepaspectratio]{enigma.jpg}\caption{An Enigma machine}\label{fig:hist:enigma}\end{subfigure}\hfill
\begin{subfigure}[b]{0.32\textwidth}\centering\includegraphics[width=\textwidth,height=38mm,keepaspectratio]{zuse_z3.jpg}\caption{Zuse Z3 (replica)}\end{subfigure}\hfill
\begin{subfigure}[b]{0.32\textwidth}\centering\includegraphics[width=\textwidth,height=38mm,keepaspectratio]{colossus.jpg}\caption{Colossus}\end{subfigure}
\caption{Machines of the 1940s.}
\end{figure}

\subsection{John von Neumann: the stored-program computer}

In 1945 the report on the EDVAC, signed by John von Neumann, described a
hardware architecture for programmable digital computers
\citep{vonneumann1945}. Its heart is one idea.

\begin{prgconcept}[The stored program]
Program and data share the same memory: \emph{a program is data}.
Instructions are numbers kept in memory, fetched and executed one after
another.
\end{prgconcept}

The architecture has four essential units (Figure~\ref{fig:hist:vn}): a
\term{control unit}, which governs the sequence of operations; an
\term{arithmetic logic unit} (ALU), with its accumulator register; a
\term{memory}; and devices for \term{input} and \term{output}.

\begin{figure}[htb]
\centering
\begin{tikzpicture}[node distance=6mm and 8mm]
  \node[blk,fill=prgsteel,text=white,minimum width=62mm,minimum height=11mm] (m) {\textbf{Memory}: program and data};
  \node[blk,fill=prgaccent,text=white,minimum width=27mm,minimum height=15mm,below left=10mm and -27mm of m] (cu) {\textbf{Control unit}};
  \node[blk,fill=prgamber,text=white,minimum width=27mm,minimum height=15mm,below right=10mm and -27mm of m] (alu) {\textbf{ALU}\\(accumulator)};
  \node[blk,fill=prggreen,text=white,minimum width=27mm,below=10mm of cu] (in) {\textbf{Input}};
  \node[blk,fill=prggreen,text=white,minimum width=27mm,below=10mm of alu] (out) {\textbf{Output}};
  \draw[<->,thick] (cu.north) -- (cu.north|-m.south);
  \draw[<->,thick] (alu.north) -- (alu.north|-m.south);
  \draw[<->,thick] (cu) -- (alu);
  \draw[->,thick] (in) -- (cu);
  \draw[->,thick] (alu) -- (out);
\end{tikzpicture}
\caption{The four units of the von Neumann architecture.}
\label{fig:hist:vn}
\end{figure}

\begin{prgmachine}
The ESP32-C6 is still a von Neumann machine, and its four units were all at
work in Chapter~\ref{ch:firstcontact}. The registers and the ALU executed
\code{addi t1, t1, 1}; the control unit decided that the next instruction was
\ins{sw}; memory held the variable \code{counter} \emph{and the instructions
themselves}, among them the number \hex{00130313}; the joystick and the
display are input and output.
\end{prgmachine}

\subsection{Eight decades of computers}

Since then each decade has made the computer smaller, cheaper and closer to
people (Table~\ref{tab:hist:decades}, Figure~\ref{fig:hist:decades}); none of
them has changed the stored-program idea.

\begin{table}[htb]
\centering\small
\caption{Eight decades of computers.}\label{tab:hist:decades}
\begin{tabular}{@{}lll@{}}
\toprule
\textbf{Decade} & \textbf{Era} & \textbf{What changed}\\
\midrule
1940s & Wartime & ENIAC (USA) and Colossus (UK), 1943--45\\
1950s & Military and government & EDVAC, 1949: the stored program\\
1960s & Mainframe computing & one machine, many users\\
1970s & Minicomputers & a computer per department\\
1980s & Personal computing & a computer per person\\
1990s & The Internet & computers connect\\
2000s & Smartphones and tablets & a computer in the pocket\\
2010s & Cloud computing & computing as a service\\
2020s & Internet of Things and AI & computers everywhere\\
\bottomrule
\end{tabular}
\end{table}

\begin{figure}[htb]
\centering
\begin{subfigure}[b]{0.24\textwidth}\centering\includegraphics[width=\textwidth,height=28mm,keepaspectratio]{eniac.png}\caption{ENIAC}\end{subfigure}\hfill
\begin{subfigure}[b]{0.24\textwidth}\centering\includegraphics[width=\textwidth,height=28mm,keepaspectratio]{mainframe.jpg}\caption{A mainframe}\end{subfigure}\hfill
\begin{subfigure}[b]{0.24\textwidth}\centering\includegraphics[width=\textwidth,height=28mm,keepaspectratio]{minicomputer.jpg}\caption{A minicomputer}\end{subfigure}\hfill
\begin{subfigure}[b]{0.24\textwidth}\centering\includegraphics[width=\textwidth,height=28mm,keepaspectratio]{c64.jpg}\caption{Commodore 64}\end{subfigure}
\caption{From a room to a desk in forty years.}
\label{fig:hist:decades}
\end{figure}

% ----------------------------------------------------------------------------
\section{From home computers to RISC-V}

\subsection{Microprocessors: from CISC to RISC to an open standard}

The processor of the \prg kit is the way it is because of a sequence of
decisions taken over fifty years (Table~\ref{tab:hist:cpu}). The first
microprocessors grew by accretion: each generation added instructions, often
complicated ones, to make the work of compilers and of assembly programmers
easier. The result is called a \term{complex instruction set computer}
(CISC), and the x86 family, born in 1978, still carries that inheritance in
today's personal computers.

At the beginning of the 1980s two university projects, at Berkeley and at
Stanford, measured what programs actually do, and found that a processor with
few, simple and regular instructions can be made to run faster than one with
many complicated ones. This is the idea of the \term{reduced instruction set
computer} (RISC), for which David Patterson and John Hennessy received the
Turing Award in 2017. ARM applied it to low power, and ARM cores are today in
nearly every telephone. In 2010 the fifth RISC design of Berkeley became
something new: an instruction set published as an \emph{open standard}.

\begin{table}[htb]
\centering\small
\caption{Microprocessors: from CISC to RISC to an open standard.}\label{tab:hist:cpu}
\begin{tabularx}{\textwidth}{@{}lP{0.31\textwidth}L@{}}
\toprule
\textbf{Year} & \textbf{Milestone} & \textbf{Why it matters}\\
\midrule
1971 & Intel 4004 & the first commercial microprocessor: 2\,300 transistors on one chip\\
1975 & MOS 6502 & a cheap 8-bit CPU: Apple II, Commodore 64, the home-computer era\\
1978 & Intel 8086 & the start of x86, a complex instruction set (CISC) still in today's PCs\\
1980--82 & Berkeley RISC and Stanford MIPS & few, simple, regular instructions run faster: the RISC idea\\
1985 & ARM1 & RISC for low power: today in nearly every telephone\\
2010 & \riscv starts at UC Berkeley & the fifth Berkeley RISC design becomes an open instruction set\\
2015--20 & RISC-V Foundation, then RISC-V International & a shared specification, governed from Switzerland\\
2023--25 & European Chips Act; the DARE project & \riscv as a strategy for European technological sovereignty\\
\bottomrule
\end{tabularx}
\end{table}

\begin{definition}{RISC-V}
\riscv is an open standard \term{instruction set architecture}: the
specification of the instructions, of the registers and of their encoding is
public, and anyone may design, build and sell a processor that implements it
without paying a licence for the architecture.
\end{definition}

\riscv is therefore not a processor, and it was not the first RISC: Berkeley
RISC, MIPS and ARM came earlier. What distinguishes it from x86 and from ARM
is openness.

\subsection{Moore's law in one number}

The Intel 4004 of 1971 had 2\,300 transistors. A laptop processor of today
has tens of thousands of millions: a growth by a factor of about ten million.
For fifty years the number of transistors on a chip has doubled about every
two years, an observation made by Gordon Moore in 1965 and revised in 1975
\citep{moore1965}. It is an observation about an industry, not a law of
nature, and its physical limits are now close. Progress comes increasingly
from better architectures, from specialised accelerators and from open
designs, which is why architecture, the subject of this book, matters more
than ever. Chapter~\ref{ch:beyond} takes up the story from here.

\subsection{The home-computer spirit, on a modern open CPU}

The home computers of the 1980s had a quality that later machines lost: one
person could understand all of one. \prg was designed to give that quality
back, on the architecture that students will meet in industry
(Table~\ref{tab:hist:c64}). The screen has the same size, and the same
immediacy, forty years apart.

\begin{table}[htb]
\centering\small
\caption{A Commodore~64 and \prg.}\label{tab:hist:c64}
\begin{tabularx}{\textwidth}{@{}lLL@{}}
\toprule
 & \textbf{Commodore 64 (1982), MOS 6510} & \textbf{\prg on ESP32-C6, \riscv}\\
\midrule
Word size & 8 bits & 32 bits\\
Registers & A, X, Y (plus stack pointer, status, PC) & 32 general-purpose registers, \reg{x0} to \reg{x31}\\
Clock & about 1 MHz & 160 MHz\\
Address space & 64 KB & 4 GB\\
Instruction set & 56 instructions, proprietary & about 40 base instructions plus extensions, open standard\\
Screen & 320 by 200 pixels, 16 colours & 320 by 200 viewport, 256-entry RGB565 palette\\
Programmed by & typing on the machine itself & building a cartridge and uploading it\\
\bottomrule
\end{tabularx}
\end{table}

The 6502 is still a pleasure to explore, for instance with the simulator at
\url{https://skilldrick.github.io/easy6502}.

\subsection{Retro games as specimens}

A biologist studies a fruit fly not because flies are important but because a
fly is small enough to be understood completely and rich enough to show how
heredity works. A retro game plays the same part here
(Figure~\ref{fig:hist:games}): it is small enough to be held in one's head and
rich enough to exercise state, input, arithmetic, decisions and a loop, which
is to say the whole course. The four games of the figure are example
cartridges of the \prg repository (\fname{examples/games}).

\begin{figure}[htb]
\centering
\begin{subfigure}[b]{0.24\textwidth}\includegraphics[width=\textwidth]{game_pong.png}\caption{Pong}\end{subfigure}\hfill
\begin{subfigure}[b]{0.24\textwidth}\includegraphics[width=\textwidth]{game_breakout.png}\caption{Breakout}\end{subfigure}\hfill
\begin{subfigure}[b]{0.24\textwidth}\includegraphics[width=\textwidth]{game_tetris.png}\caption{Tetris}\end{subfigure}\hfill
\begin{subfigure}[b]{0.24\textwidth}\includegraphics[width=\textwidth]{game_asteroids.png}\caption{Asteroids}\end{subfigure}
\caption{Retro games as specimens.}
\label{fig:hist:games}
\end{figure}

% ----------------------------------------------------------------------------
\section{Informatics as an applied science}

Where does informatics sit among the sciences? An example from the
Department in which this book was written makes the answer concrete
(Figure~\ref{fig:hist:chain}). A storm forms over the sea: a natural
phenomenon. Physics and mathematics, the basic sciences, describe it with
equations. Informatics, an applied science, turns the equations into
algorithms and simulations. Engineering builds the supercomputers, the
networks and the sensors that run them reliably. And the economy receives
forecasts and services that reach people and create value. The weather and
marine forecasts of \url{https://meteo.uniparthenope.it} are computed every
day in this way, on the Department's own high-performance systems.

\begin{figure}[htb]
\centering
\begin{tikzpicture}[node distance=4mm]
  \node[blkacc,text width=22mm,minimum height=17mm] (a) {\textbf{Natural phenomenon}\\[1pt]\footnotesize a storm};
  \node[blkamb,text width=22mm,minimum height=17mm,right=of a] (b) {\textbf{Basic sciences}\\[1pt]\footnotesize equations};
  \node[blkvio,text width=22mm,minimum height=17mm,right=of b] (c) {\textbf{Applied sciences}\\[1pt]\footnotesize algorithms};
  \node[blkgrn,text width=22mm,minimum height=17mm,right=of c] (d) {\textbf{Engineering}\\[1pt]\footnotesize systems};
  \node[blkred,text width=22mm,minimum height=17mm,right=of d] (e) {\textbf{Economy}\\[1pt]\footnotesize services};
  \foreach \x/\y in {a/b,b/c,c/d,d/e} \draw[->,thick] (\x) -- (\y);
\end{tikzpicture}
\caption{Where informatics sits.}
\label{fig:hist:chain}
\end{figure}

The same ideas that are practised in this book on a microcontroller at
160~MHz scale up to clusters (Figure~\ref{fig:hist:hpc}): registers, memory,
instructions, only many more of them.

\begin{figure}[htb]
\centering
\begin{subfigure}[b]{0.32\textwidth}\centering\includegraphics[width=\textwidth,height=36mm,keepaspectratio]{rcf_cluster.jpg}\caption{The racks}\end{subfigure}\hfill
\begin{subfigure}[b]{0.32\textwidth}\centering\includegraphics[width=\textwidth,height=36mm,keepaspectratio]{cluster_1.jpg}\caption{Weather forecasts}\end{subfigure}\hfill
\begin{subfigure}[b]{0.32\textwidth}\centering\includegraphics[width=\textwidth,height=36mm,keepaspectratio]{rcf_page.png}\caption{\url{rcf.uniparthenope.it}}\end{subfigure}
\caption{High-performance computing at Parthenope: the racks of the
Research Computing Facilities, the forecasts of meteo@uniparthenope that they
compute every day, and the page that lists the machines.}
\label{fig:hist:hpc}
\end{figure}

% ----------------------------------------------------------------------------
% Generated by tools/textbook/export_listings.js from decks/plan.js and code/l03.S: do not edit.
\section{The program of this chapter}\label{sec:lab:l03}

The listings of this chapter are parts of one program, \fname{code/l03.S} in the
repository of the course, printed in full in Listing~\ref{lst:lab:l03}. The
automatic tests of the course material build it and run it on the headless
runner of \prgqt{}, with scripted input, and check the frame;
Figure~\ref{fig:lab:l03} is the frame that the test wrote.

\begin{description}[style=nextline,leftmargin=1.5em]
\item[On the screen.] A row of 32 lamps showing 0x00130313, the machine code of addi t1, t1, 1, and a green square: the word read is the one expected.
\item[What to press.] Nothing: the program runs by itself.
\end{description}

\begin{figure}[htb]
\centering
\includegraphics[width=0.5\textwidth]{lab_l03.png}
\caption{\fname{l03.S} on \prgqt{}.}\label{fig:lab:l03}
\end{figure}

\asmfile[caption={The complete program of the chapter, \fname{code/l03.S}.},label=lst:lab:l03]{code/l03/program.S}

To run it on \prgqt{} the commands are those of Listing~\ref{lst:lab:l03:run};
Section~\ref{sec:prg:testing} gives the same steps for QEMU and for the board.

\begin{shell}[caption={Building \fname{l03.S} and running it on \prgqt{}.},label=lst:lab:l03:run]
# 1. start PRG32-QT on this computer (it listens on port 8080)

# 2. in the PRG32 folder, with ESP-IDF active: build the cartridge
python3 -m prg32 cartridge build ~/arc12/code/l03.S --portable \
    --entry-prefix l03 --name l03 --out l03.prg32

# 3. send it to PRG32-QT and start it
python3 -m prg32 esp32c6 upload-and-run l03.prg32 \
    --url http://localhost:8080 --slot cart0

# 4. look, press, compare with the previous slide
\end{shell}


% ----------------------------------------------------------------------------
% Generated by tools/textbook/export_listings.js from tools/concepts/assessments.json: do not edit.
\begin{chapstatus}{The thread of this chapter}
The machine you programmed in Lecture 02 is a von Neumann machine: program and data in one memory. RISC-V is the latest step of a long debate on how simple instructions should be. Next: how anything at all becomes bits.
\par\medskip\textbf{Exit ticket} (two to four minutes, alone, on paper):
\begin{enumerate}[nosep,leftmargin=2em]
\item Explain the sentence attributed to Dijkstra, 'Computer science is no more about computers than astronomy is about telescopes', using the definition of informatics given in the lecture.
\item Where is the program of a cartridge kept while it runs, and why does this make a von Neumann machine different from a machine programmed by rewiring cables?
\end{enumerate}
\end{chapstatus}


\section{Check your understanding}

\begin{quiz}{Which idea is the heart of the von Neumann architecture?}
\item The processor has exactly four registers
\item Program and data are stored in the same memory
\item Every instruction takes one clock cycle
\item Input and output are optional
\end{quiz}

\begin{quiz}{What makes \riscv different from x86 and ARM?}
\item It is the fastest instruction set
\item It was the first RISC architecture
\item Its instruction set is an open standard, free to implement
\item It works only on microcontrollers
\end{quiz}

\begin{quiz}{Who is regarded as the first programmer in history?}
\item Charles Babbage
\item Alan Turing
\item Ada Lovelace
\item Konrad Zuse
\end{quiz}

\section{Exercises}

\begin{exercise}[paper first, then \prg]
\begin{enumerate}[nosep]
\item Draw the four units of the von Neumann architecture and write, inside
  each, one thing of the program of Chapter~\ref{ch:firstcontact} that lives
  or happens there.
\item For each of the following, say whether it is a new way of
  \emph{representing} information or a new way of \emph{prescribing}
  operations: the abacus, the Jacquard loom, the Pascaline, Boole's algebra,
  the punched cards of the Analytical Engine, the stored program.
\item Download two of the four games of Figure~\ref{fig:hist:games} from the
  Store or build them from \fname{examples/games}. With the counters of the
  debugger, find how many instructions each executes per frame, and compare
  the two games.
\end{enumerate}
\expected{Registers and ALU: \ins{addi} was executed there. Memory: the
counter, and the instructions themselves. Control unit: the fetch of each
instruction. Input/output: the bar on the screen. The stored program belongs
to both classes: a program is itself represented information.}
\end{exercise}

\section{Glossary}

\begin{glossarytable}
computer science / informatics & informatica\\
computer & calcolatore, elaboratore\\
punched card & scheda perforata\\
loom & telaio\\
mainframe & mainframe, grande elaboratore\\
stored program & programma memorizzato\\
control unit & unit\`a di controllo\\
arithmetic logic unit & unit\`a aritmetico-logica\\
instruction set & insieme (set) di istruzioni\\
high-performance computing & calcolo ad alte prestazioni\\
\end{glossarytable}

\section{Conclusions and reading}

Computing is older than electronics: representation and procedure came first.
The stored-program computer of 1945 is still the machine that is programmed
today. RISC made processors simple enough to be understood whole; \riscv made
them open. Informatics is pervasive: digital information belongs to every
field, data are its currency, and algorithms turn data into decisions.

\paragraph{Reading.} \citet{patterson2020}, chapter~1; \emph{Playing with
RISC-V}, chapter~1, ``Why a Game, Why RISC-V, Why Now''
\citep{prg32textbook}; Charles Petzold, \emph{Code: The Hidden Language of
Computer Hardware and Software}; for the history in depth,
\citet{ceruzzi2003}.
